\documentclass[11pt]{article}
\newcommand{\shorttitle}{Prime sums and zeta zeros}
\usepackage[T1]{fontenc}
\usepackage{lmodern,amsmath,amssymb,amsthm,mathtools,booktabs,microtype}
\usepackage[letterpaper,left=1.03in,right=1.03in,top=.90in,bottom=.90in]{geometry}
\usepackage{xcolor,fancyhdr,enumitem}
\usepackage{hyperref}
\definecolor{accent}{HTML}{1D5660}
\definecolor{ink}{HTML}{192E38}
\hypersetup{colorlinks=true,linkcolor=accent,urlcolor=accent,citecolor=accent}
\newtheorem{theorem}{Theorem}[section]
\newtheorem{corollary}[theorem]{Corollary}
\newtheorem{proposition}[theorem]{Proposition}
\setcounter{tocdepth}{1}
\setlength{\headheight}{14pt}
\setlength{\emergencystretch}{1.3em}
\setlength{\parindent}{1.15em}
\setlist{itemsep=0.15em,topsep=0.4em}
\pagestyle{fancy}\fancyhf{}
\fancyhead[L]{\small\color{accent}\shorttitle}
\fancyhead[R]{\small\color{accent}25 November 2026}
\fancyfoot[C]{\small\thepage}
\renewcommand{\headrulewidth}{0.3pt}
\newcommand{\paperheading}[3]{\noindent{\sffamily\small\bfseries\color{accent}#1}\par\vspace{.75em}
  {\raggedright\sffamily\bfseries\LARGE\color{ink}#2\par}\vspace{.6em}
  {\noindent\large #3\par}\vspace{1em}}

\hypersetup{pdftitle={What sparse prime sums determine about zeta zeros}}
\begin{document}
\paperheading{ANALYTIC NUMBER THEORY}{What sparse prime sums determine about zeta zeros}{Two transfers of the one-sided prime-remainder method}
\noindent
For every fixed integer $k\ge1$, a one-sided bound on a smoothed prime-ideal sum at $x=m^k$ is equivalent to the Riemann hypothesis for the Dedekind zeta function of a given number field. The implication holds even if there is no zero with largest real part. For the ordinary Riemann zeta function, a different sequence of cutoffs, $x_n=2^n3^{n^2}$, recovers the multiplicity of every critical zero. We obtain an explicit nonnegative defect that vanishes precisely when all zeros are simple and bound it by $1.67\times10^{-34}$. The infinite estimate for the nonabelian cubic-field example and the vanishing of the simplicity defect remain open. Both questions require information beyond the five conclusions of the Lean repository.
\section{A one-sided test on sparse scales for a Dedekind zeta function}

Let $K$ be a number field. Write $\Lambda_K(n)$ for the coefficient of
$n^{-s}$ in $-\zeta'_K/\zeta_K(s)$, and $\gamma_K$ for the constant term
of $\zeta'_K/\zeta_K(s)$ at $s=1$. The analogue of the repository's ramp is
\[
 R_K(x)=\sum_{n\le x}\Lambda_K(n)(1/n-1/x)-\log x+1+\gamma_K ,
 \qquad x\ge1.
\]
It is continuous at a prime-power norm because its entering summand vanishes.
Let $\Theta_K=\sup_\rho\Re\rho$, over the nontrivial zeros of $\zeta_K$,
counting multiplicity. The functional equation gives $1/2\le\Theta_K\le1$.

\paragraph{Theorem (prescribed powers).}
Fix $k\ge1$. For every $0<\beta<\Theta_K$,
\[
 \limsup_{m\to\infty}(m^k)^{1-\beta}R_K(m^k)=+\infty.
\]
For every $\beta>\Theta_K$ the same expression tends to zero.
In particular the field's Riemann hypothesis is equivalent to an eventual
one-sided bound $R_K(m^k)\le C_{K,k}m^{-k/2}$. The result also holds on every unbounded increasing
grid $y_j$ with $\log(y_{j+1}/y_j)=O(y_j^{-\delta})$ for some $\delta>0$.

\paragraph{Proof.}
For $\Re s>1$ the Mellin transform $H(s)=\int_1^\infty R_K(x)x^{-s}\,dx$ is
\[
 H(s)=\frac{-\zeta'_K(s)/\zeta_K(s)}{s(s-1)}
       -\frac{1}{(s-1)^2}+\frac{1+\gamma_K}{s-1}.
\]
The singularity at one is removable. At each nontrivial zero $\rho$ its
residue is $m_\rho/[\rho(1-\rho)]$, strictly positive at real zeros.
Consequently, an eventual bound $R_K(x)\le Cx^{\tau-1}$, $0<\tau<1$,
excludes all zeros in $\Re s>\tau$. Indeed, after increasing $C$ the
Mellin transform of $Cx^{\tau-1}-R_K(x)\ge0$ is $C/(s-\tau)-H(s)$.
If its real abscissa of convergence exceeded $\tau$, Landau's theorem
would give a real singularity there. Every possible pole there is simple
with negative residue, incompatible with nonnegative values approaching
from the right. Thus the transform is holomorphic on $\Re s>\tau$.
This uses the sign argument formalized in
\href{https://github.com/lsernau/grh-lean/blob/72fead102f663d5c7a407964a13dc0cf9d977207/grh/src/Lib/Analytic/Landau.lean}{\texttt{IsPolar.landau}}.

The smoothed explicit formula is
\[
 R_K(x)=\sum_\rho\frac{x^{\rho-1}}{\rho(1-\rho)}
             +O_K(x^{-1}\log(2x)),\qquad x\ge2 .
\]
Here and below each zero is repeated according to multiplicity.
One obtains it by shifting the inverse Mellin integral for the uncentered
ramp from $\Re s=2$ to $-1/2$, taking heights avoiding zeros. The terms
at one give the displayed centering, the at-most-double pole at zero gives
$O_K(x^{-1}\log(2x))$, and the new vertical line contributes
$O_K(x^{-3/2})$ by the functional equation and Stirling's formula.
Absolute convergence follows from the zero count $N_K(T)=O_K(T\log(2T))$
for $T\ge2$. Compare Y.~Ihara,
\href{https://www.kurims.kyoto-u.ac.jp/~kenkyubu/emeritus/ihara/Publications-and-Recent-Preprints/RecentArticles/pdf-files/1A.pdf}
{\emph{On the Euler--Kronecker constants of global fields and primes with small norms}}, Part~1.

For $0<h\le1$ one has $|e^{(\rho-1)h}-1|\le\min(2,|\rho-1|h)$.
Splitting the sum at $|\Im\rho|=1/h$ gives
\[
 |R_K(xe^h)-R_K(x)|\ll_K
 x^{\Theta_K-1}h\log^2(2/h)+x^{-1}\log(2x).
\]
Indeed partial summation of the zero count gives
$\sum_{|\Im\rho|\le T}|\rho|^{-1}\ll_K\log^2(2T)$ and
$\sum_{|\Im\rho|>T}|\rho(1-\rho)|^{-1}\ll_K T^{-1}\log(2T)$ for $T\ge1$.
If the asserted limsup at $\beta<\Theta_K$ were bounded above,
the grid values would satisfy $R_K(y_j)\ll y_j^{\beta-1}$ in the one-sided
sense. Choose $\max(\beta,\Theta_K-\delta)<\tau<\Theta_K$.
For $y_j\le x\le y_{j+1}$ the variation estimate yields
$R_K(x)\le Cx^{\tau-1}$ throughout the ray. The Mellin argument then
contradicts the definition of the supremum. For $y_m=m^k$ one takes
$\delta=1/k$. For $\beta>\Theta_K$ the conclusion follows directly by
absolute convergence of the explicit formula. On the critical line the same argument gives $|R_K(x)|\ll_K x^{-1/2}$, including the exact endpoint exponent.\hfill$\square$

\paragraph{An equivalent finite statistic.}
The preceding criterion can be evaluated without $\gamma_K$. Put
$D_k(m)=R_K(m^k)-R_K((2m)^k)$, obtained by subtracting the two finite ramp sums and adding $k\log2$. Then RH for $K$ is equivalent to $\sup_{m\ge1}m^{k/2}D_k(m)<\infty$. On RH the two-sided endpoint bound proves this inequality. Conversely, $R_K(x)\to0$ unconditionally by the explicit formula, dominated convergence, and nonvanishing on $\Re s=1$. Telescoping over $m,2m,4m,\ldots$ gives $R_K(m^k)\le C m^{-k/2}/(1-2^{-k/2})$ from a nonnegative upper constant $C$ for the statistic. The statistic extends the quartic-grid test of the incoming report \texttt{grh\_method\_transfers\_2026-11-25.zip} to every prescribed power for the full prime-ideal weights. That report also gives a prime-only version at fourth powers.

\paragraph{A single base for every number field.}
Let $0<\lambda_n\uparrow\infty$ with $\lambda_{n+1}/\lambda_n\to1$.
There is a dense countable intersection of open sets of bases $b>1$
on which the same two limsup conclusions hold at $x=b^{\lambda_n}$
for every $K$ simultaneously. For instance,
$\lambda_n=\exp(n/\log(n+2))$ needs asymptotically
$(\log\log X)\log\log\log X$ samples up to $X$ for fixed $b$.
To see this, for any continuous $F$ and rational $h<\limsup_{x\to\infty}F(x)$
the set of $b$ for which $F(b^{\lambda_n})>h$ for some $n\ge N$ is open
and dense. Indeed for $1<a<c$ the intervals
$(a^{\lambda_n},c^{\lambda_n})$ eventually overlap and cover a ray.
Intersect over $h,N$ and the countably many functions
$F(x)=x^{1-\beta}R_K(x)$ for rational $0<\beta<1$. Landau's argument gives
full limsup $+\infty$ for $\beta<\Theta_K$; intermediate rational exponents
give the real exponents. This is a Croft--Kingman sampling argument.
\href{https://doi.org/10.1017/CBO9781139107174.008}{Bingham and Ostaszewski,
\emph{Kingman, category and combinatorics}} give its topological background.
The generic-base theorem leaves the status of a prescribed base such as $2$
open on these much thinner scales.

\paragraph{A nonabelian example.}
Take $K=\mathbb Q(\sqrt[3]{2})$, of discriminant $-108$ and Galois closure
group $S_3$. For $p\nmid6$ the degrees $f_i$ of the irreducible factors of
$T^3-2$ modulo $p$ give $\Lambda_K(p^j)=(\log p)\sum_{f_i\mid j}f_i$.
At $2$ the polynomial $T^3-2$ is Eisenstein, and at $3$ the translate $(T-1)^3-2$ is Eisenstein. At either prime there is one totally ramified prime of residue degree one, so
$\Lambda_K(p^j)=\log p$. Thus the prescribed-power test has computable
prime weights even here. The permutation representation of $S_3$ gives
$\zeta_K=\zeta\,L(s,\mathrm{std})$; the two-dimensional standard factor
is the input beyond Dirichlet GRH. The criterion supplies an exact
arithmetic target. Its infinite sampled upper bound for this field
remains to be proved.\quad
\href{https://www.lmfdb.org/NumberField/3.1.108.1}{LMFDB field 3.1.108.1}.

\section{A sparse prime-sum criterion for simple zeros}
\label{sec:sparse-simplicity}

The logarithmic-derivative method behind the Lean reduction also detects
multiplicity. For the Riemann zeta function it gives the following precise
target. Let $\gamma$ denote Euler's constant and $\gamma_1$ the first Stieltjes
constant, with $\zeta(1+u)=u^{-1}+\gamma-\gamma_1u+O(u^2)$. Set
\[
 x_n=2^n3^{n^2},\qquad
 R_{\mathbb Q}(x)=
 \sum_{p^j\le x,\ j\ge1}\log p\left(\frac{1}{p^j}-\frac1x\right)
 -\log x+1+\gamma .
\]
\noindent\textbf{Simplicity criterion.}
Under RH, all nontrivial zeros of $\zeta$ are simple if and only if
\[
 \lim_{N\to\infty}\frac1N\sum_{n=1}^N x_nR_{\mathbb Q}(x_n)^2
 =6+2\gamma-2\log(4\pi)+4\gamma_1+2\gamma^2-\frac{\pi^2}{4}.
\]
The constant is approximately $0.00007420127287493$. The limit always
exists under RH. Its excess above the displayed value is positive if any
zero is multiple. The remaining arithmetic question is an upper bound on
the sampled mean square by this explicit constant. Only
$O(\sqrt{\log X})$ cutoffs occur through $X$, although each cutoff
uses prime powers up to its full value. Continuous mean squares of prime
errors also detect repeated zeros. Here the reciprocal-prime-power
weight makes the Fourier series absolutely convergent, and the polynomial
choice of logarithms permits sampling on a single sparse sequence.
\footnote{For the continuous mean square of the unsmoothed Chebyshev error
and its dependence on multiplicities, see P.~Caceres,
\emph{The constant term in Cramér's mean-square theorem for the primes},
\href{https://hexagonmath.org/2611.01010v1}{preprint, 2026}.}

\paragraph{The zero expansion on the critical line.}
Use the same $R_K$ and explicit formula as in Section~1. In Ihara's notation
the uncentered prime sum equals $(1-x^{-1})\Phi_K(x)$.
The normalization and archive sources are recorded in \texttt{sources.md}.

Assume all nontrivial zeros of $\zeta_K$ lie on the critical line. For each
real $t$ let $m_K(t)$ be the multiplicity of $1/2+it$, with value zero at
nonzeros, and let $\Gamma_K$ be the set of distinct ordinates. Then
\[
 \sqrt{x}R_K(x)=\sum_{t\in\Gamma_K} a_t x^{it}+o(1),
 \qquad a_t=\frac{m_K(t)}{1/4+t^2},\qquad \sum_t a_t<\infty .
\]
Absolute summability gives uniform convergence in $\log x$. The Lean
repository supplies this line hypothesis when $K=\mathbb Q$. For abelian
extensions of $\mathbb Q$ it follows from the factorization into Dirichlet
$L$-functions. For a general number field it is an additional hypothesis.

\paragraph{Recovery of every multiplicity.}
For every real $u$,
\[
 \lim_{N\to\infty}\frac1N\sum_{n=1}^N x_n^{iu}
   =\begin{cases}1&u=0,\\0&u\ne0.\end{cases}
\]
Indeed, for $u\ne0$, at least one of $u\log2/(2\pi)$ and $u\log3/(2\pi)$ is
irrational, since $\log2/\log3$ is irrational by unique factorization. Apply
Weyl's theorem to the polynomial with these linear and quadratic
coefficients. The same argument works with $n^2$ replaced by $n^d$ for
each fixed integer $d\ge2$.
\footnote{H.~Weyl, \emph{\"Uber die Gleichverteilung von Zahlen mod.\ Eins},
Math. Ann. 77 (1916), 313--352, \href{https://doi.org/10.1007/BF01475864}{doi:10.1007/BF01475864}.}
Multiplication by $x_n^{-iu}$ and termwise averaging now give
\[
 \lim_{N\to\infty}\frac1N\sum_{n=1}^N
 \sqrt{x_n}R_K(x_n)x_n^{-iu}=\frac{m_K(u)}{1/4+u^2}.
\]
This recovery formula requires no assumptions about rational relations
among zero ordinates. Absolute convergence controls the tail uniformly in
$N$, and the average of the $o(1)$ term tends to zero.

\paragraph{The defect from simplicity.}
Since the prime sum is real, its square is its squared absolute value.
Applying the same averaging identity to pairs of ordinates gives
\[
 \lim_{N\to\infty}\frac1N\sum_{n=1}^N x_nR_K(x_n)^2
     =\sum_{t\in\Gamma_K}\frac{m_K(t)^2}{(1/4+t^2)^2}.
\]
Indeed the double sum is absolutely convergent, and only terms whose
ordinates coincide survive. Write $H_K=(\log\xi_K)'$, where
$\xi_K(s)=s(s-1)|D_K|^{s/2}\Gamma_{\mathbb R}(s)^{r_1}
\Gamma_{\mathbb C}(s)^{r_2}\zeta_K(s)$ is the completed entire zeta
function, $D_K$ is the discriminant and $(r_1,r_2)$ the signature.
Here $\Gamma_{\mathbb R}(s)=\pi^{-s/2}\Gamma(s/2)$ and
$\Gamma_{\mathbb C}(s)=2(2\pi)^{-s}\Gamma(s)$.
The genus-one Hadamard product gives
\[
 \sum_\rho\frac1{\rho(1-\rho)}
       =H_K(1)-H_K(0)=2H_K(1),\qquad
 \sum_\rho\frac1{(1-\rho)^2}=-H'_K(1).
\]
The first identity uses $\xi_K(s)=\xi_K(1-s)$.
The identity
$[z(1-z)]^{-2}=z^{-2}+(1-z)^{-2}+2/[z(1-z)]$ and the same symmetry give
\[
 \sum_{t\in\Gamma_K}\frac{m_K(t)}{(1/4+t^2)^2}
      =4H_K(1)-2H'_K(1).
\]
Consequently the difference between the sampled mean square and
$4H_K(1)-2H'_K(1)$ is exactly
$\sum_t m_K(t)(m_K(t)-1)/(1/4+t^2)^2$. Every term is nonnegative.
For $K=\mathbb Q$ the expansion at $s=1$ gives
\[
 H_{\mathbb Q}(1)=1+\tfrac12\gamma-\tfrac12\log(4\pi),\qquad
 H'_{\mathbb Q}(1)=-1-2\gamma_1-\gamma^2+\pi^2/8,
\]
which proves the stated constant. Controlling the off-diagonal terms when
one expands the \emph{prime} side of the mean square at these cutoffs
remains open.

\paragraph{A numerical enclosure for the limiting statistic.}
Let $D$ be the excess in the Riemann case and $H=3\cdot10^{12}$.
Then
\[
 0\le D <1.67\cdot10^{-34}.
\]
Platt and Trudgian's disjoint sign changes exhaust the zero count with multiplicity through $H$, so each of those zeros is simple. For $t\ge H$ Corollary~1.2 of Hasanalizade, Shen and Wong
bounds the difference of the zero count $N(t)$ from
$(t/2\pi)\log(t/(2\pi e))$ by
$0.1038\log t+0.2573\log\log t+9.3675<\log t$.
Hence $N(t)\le t\log t$ there, and a jump at $t>H$ has size at most
$2\log t$. By conjugation symmetry and Stieltjes integration by parts,
\[
 D\le 4\int_{(H,\infty)}\frac{\log t}{t^4}\,dN(t)
 \le16\int_H^\infty\frac{\log^2t}{t^4}\,dt
 =\frac{16}{H^3}\left(\frac{\log^2H}3+
     \frac{2\log H}9+\frac2{27}\right).
\]
The boundary term discarded at $H$ is nonpositive. A multiple zero at
arbitrarily large height could contribute less than this bound.
\footnote{D.~Platt and T.~Trudgian,
\emph{The Riemann hypothesis is true up to $3\cdot10^{12}$},
\href{https://doi.org/10.1112/blms.12460}{Bull. London Math. Soc. 53 (2021), 792--797}.
E.~Hasanalizade, Q.~Shen and P.-J.~Wong,
\href{https://arxiv.org/abs/2107.06506}{\emph{Counting zeros of the Riemann zeta function}},
Corollary~1.2. The earlier ProveIt report on Suzuki's density question uses
the same two estimates with a different spectral weight.}

\end{document}
